\section{Conclusion}
\label{sec:conclusion}

We studied how uncertainty estimation and deferral policies interact in
a medical image segmentation pipeline, with the goal of making reliable
\emph{decisions} rather than producing better segmentation maps. Three
findings stand out.

First, the choice of uncertainty method matters more for deferral than for
segmentation. TTA and MC Dropout produce similar Dice scores (0.768 vs.\
0.763), but TTA's uncertainty is far more useful for identifying errors
(Unc-AUROC 0.881 vs.\ 0.722 in our experiments). If the purpose of uncertainty is to support
deferral, TTA is the better choice on both quality and speed.

Second, how uncertainty is used matters as much as how it is generated.
TTA with adaptive deferral reduces segmentation error by approximately 80\%
at a 25\% deferral rate. Confidence-aware deferral offers the best efficiency:
55\% error reduction at just 12\% deferral. Both far exceed the weakest
operating points in the same comparison family. The deferral policy is not an afterthought;
it is half the system. Uncertainty without a decision mechanism to act on it
leaves most of this error reduction unrealized.

Third, calibration is not a prerequisite for good deferral. Temperature
scaling does not improve the operating points we tracked and
it worsens the reported ECE values in our experiments.
Calibration optimizes a distributional property (probabilities match
frequencies in aggregate) while deferral requires a discriminative one
(uncertainty separates correct from incorrect pixels). These objectives
are related, but they are not interchangeable.

Future work should extend this analysis to multi-class segmentation and
3D imaging, where uncertainty aggregation and deferral at the voxel, slice,
or volume level introduce additional design choices. Integrating
human-in-the-loop evaluation, where clinicians interact with deferral maps
and their diagnostic accuracy is measured, would close the gap between
computational deferral quality and clinical utility.
